\documentclass[12pt,a4paper]{article}

\usepackage[utf8]{inputenc}
\usepackage[vietnamese]{babel}
\usepackage{amsmath,amssymb,amsthm}
\usepackage{geometry}
\usepackage{enumitem}
\usepackage[most]{tcolorbox}
\usepackage{hyperref}
\usepackage{fancyhdr}
\usepackage{tikz}
\usepackage{eso-pic}
\usepackage{xcolor}

\geometry{a4paper, margin=2cm, bottom=2.5cm}
\hypersetup{colorlinks=true, linkcolor=blue!70!black}
\setlist[itemize]{topsep=4pt, itemsep=2pt}
\setlist[enumerate]{topsep=4pt, itemsep=2pt}

\AddToShipoutPictureFG{%
  \AtPageCenter{%
    \begin{tikzpicture}[remember picture, overlay]
      \node[rotate=45, scale=1, opacity=0.06]{\fontsize{60}{72}\selectfont\bfseries\sffamily Đặng Tấn Quí};
    \end{tikzpicture}%
  }%
}

\pagestyle{fancy}
\fancyhf{}
\fancyhead[L]{\small\textit{Đại số cơ sở}}
\fancyhead[R]{\small\textit{Đặng Tấn Quí}}
\fancyfoot[C]{\thepage}
\fancyfoot[L]{\footnotesize\textcopyright\ Đặng Tấn Quí}
\fancyfoot[R]{\footnotesize SĐT: 0377\,122\,210}
\renewcommand{\headrulewidth}{0.4pt}
\renewcommand{\footrulewidth}{0.4pt}

\fancypagestyle{plain}{\fancyhf{}\fancyfoot[C]{\thepage}\fancyfoot[L]{\footnotesize\textcopyright\ Đặng Tấn Quí}\fancyfoot[R]{\footnotesize SĐT: 0377\,122\,210}\renewcommand{\headrulewidth}{0pt}\renewcommand{\footrulewidth}{0.4pt}}

\newtcolorbox{dinhnghia}[1][]{enhanced, breakable, colback=blue!3!white, colframe=blue!60!black, fonttitle=\bfseries, title={Định nghĩa}, #1}
\newtcolorbox{dinhly}[1][]{enhanced, breakable, colback=green!3!white, colframe=green!50!black, fonttitle=\bfseries, title={Định lý}, #1}
\newtcolorbox{tinhchat}[1][]{enhanced, breakable, colback=yellow!5!white, colframe=orange!50!black, fonttitle=\bfseries, title={Tính chất}, #1}
\newtcolorbox{phuongphap}[1][]{enhanced, breakable, colback=gray!5!white, colframe=gray!50!black, fonttitle=\bfseries, title={Phương pháp}, #1}
\newtcolorbox{luuy}[1][]{enhanced, breakable, colback=red!3!white, colframe=red!50!black, fonttitle=\bfseries, title={Lưu ý}, #1}
\newtcolorbox{congthuc}[1][]{enhanced, breakable, colback=cyan!3!white, colframe=cyan!50!black, fonttitle=\bfseries, title={Công thức}, #1}
\newtcolorbox{vidu}[1][]{enhanced, breakable, colback=violet!3!white, colframe=violet!50!black, fonttitle=\bfseries, title={Ví dụ}, #1}

\begin{document}

\thispagestyle{empty}
\begin{tikzpicture}[remember picture, overlay]
  \fill[blue!80!black] (current page.south west) rectangle (current page.north east);
  \node[anchor=west, white, font=\fontsize{26}{32}\bfseries\selectfont]
    at ([xshift=3cm, yshift=-7.5cm]current page.north west) {ĐẠI SỐ CƠ SỞ};
  \node[anchor=west, white, font=\fontsize{18}{24}\selectfont]
    at ([xshift=3cm, yshift=-9cm]current page.north west) {Tập hợp, ánh xạ, số phức --- Năm 1, HK1};
  \node[anchor=west, white, font=\Large\bfseries] at ([xshift=3cm, yshift=-13cm]current page.north west) {Biên soạn:};
  \node[anchor=west, yellow!80!white, font=\fontsize{22}{28}\bfseries\selectfont] at ([xshift=3cm, yshift=-14.5cm]current page.north west) {ĐẶNG TẤN QUÍ};
  \node[anchor=south east, white, font=\Large\bfseries] at ([xshift=-2cm, yshift=3cm]current page.south east) {\the\year};
\end{tikzpicture}

\newpage
\tableofcontents
\newpage

\section{Tập hợp}

\subsection{Các khái niệm cơ bản}

\begin{dinhnghia}
  \textbf{Tập hợp}: tập các đối tượng phân biệt. $x \in A$: $x$ thuộc $A$; $x \notin A$: $x$ không thuộc $A$.
\end{dinhnghia}

\begin{dinhnghia}
  \textbf{Tập con}: $A \subset B$ $\Leftrightarrow$ $\forall x \in A \Rightarrow x \in B$. \textbf{Bằng nhau}: $A = B$ $\Leftrightarrow$ $A \subset B$ và $B \subset A$.
\end{dinhnghia}

\begin{dinhnghia}
  \textbf{Tập rỗng} $\varnothing$: tập không chứa phần tử nào. $\varnothing \subset A$ với mọi $A$.
\end{dinhnghia}

\subsection{Các phép toán}

\begin{dinhnghia}
  \textbf{Hợp}: $A \cup B = \{x \mid x \in A \text{ hoặc } x \in B\}$.
  
  \textbf{Giao}: $A \cap B = \{x \mid x \in A \text{ và } x \in B\}$.
  
  \textbf{Hiệu}: $A \setminus B = \{x \in A \mid x \notin B\}$.
  
  \textbf{Phần bù}: $A^c = \mathcal{U} \setminus A$ (trong không gian $\mathcal{U}$).
\end{dinhnghia}

\begin{tinhchat}
  Luật De Morgan: $(A \cup B)^c = A^c \cap B^c$, $(A \cap B)^c = A^c \cup B^c$.
\end{tinhchat}

\section{Ánh xạ}

\begin{dinhnghia}
  \textbf{Ánh xạ} $f: A \to B$: quy tắc cho tương ứng mỗi $x \in A$ với đúng một $f(x) \in B$. $A$: tập nguồn, $B$: tập đích.
\end{dinhnghia}

\begin{dinhnghia}
  $f$ \textbf{đơn ánh} nếu $f(x_1) = f(x_2) \Rightarrow x_1 = x_2$.
  
  $f$ \textbf{toàn ánh} nếu $\forall y \in B$, $\exists x \in A$: $f(x) = y$.
  
  $f$ \textbf{song ánh} nếu vừa đơn ánh vừa toàn ánh.
\end{dinhnghia}

\begin{dinhnghia}
  \textbf{Ánh xạ ngược}: Nếu $f: A \to B$ song ánh thì tồn tại $f^{-1}: B \to A$ sao cho $f^{-1}(f(x)) = x$, $f(f^{-1}(y)) = y$.
\end{dinhnghia}

\section{Số phức}

\begin{dinhnghia}
  \textbf{Số phức} $z = a + bi$, $a, b \in \mathbb{R}$, $i^2 = -1$. $a$: phần thực, $b$: phần ảo.
\end{dinhnghia}

\begin{congthuc}[title={Dạng lượng giác}]
  $z = r(\cos\varphi + i\sin\varphi)$, $r = |z| = \sqrt{a^2 + b^2}$, $\varphi = \arg z$.
\end{congthuc}

\begin{congthuc}[title={Công thức Moivre}]
  $z^n = r^n(\cos n\varphi + i\sin n\varphi)$.
\end{congthuc}

\begin{dinhnghia}
  \textbf{Căn bậc $n$} của $z = r(\cos\varphi + i\sin\varphi)$:
  \[
    w_k = \sqrt[n]{r}\left(\cos\frac{\varphi + 2k\pi}{n} + i\sin\frac{\varphi + 2k\pi}{n}\right), \quad k = 0, 1, \ldots, n-1.
  \]
\end{dinhnghia}

\section{Cấu trúc đại số}

\begin{dinhnghia}
  \textbf{Nhóm} $(G, *)$: tập $G$ với phép toán $*$ thỏa kết hợp, có phần tử đơn vị, mọi phần tử có nghịch đảo.
\end{dinhnghia}

\begin{dinhnghia}
  \textbf{Vành} $(R, +, \cdot)$: $(R, +)$ là nhóm Abel, phép nhân kết hợp, phân phối với phép cộng.
\end{dinhnghia}

\begin{dinhnghia}
  \textbf{Trường}: vành giao hoán có đơn vị, mọi phần tử khác 0 có nghịch đảo. Ví dụ: $\mathbb{Q}$, $\mathbb{R}$, $\mathbb{C}$.
\end{dinhnghia}

\end{document}
